pdfoutput=1
% Options for packages loaded elsewhere
\PassOptionsToPackage{unicode}{hyperref}
\PassOptionsToPackage{hyphens}{url}
\pdfoutput=1
\documentclass[
  12pt,
]{article}
\usepackage{xcolor}
\usepackage{amsmath,amssymb}
\setcounter{secnumdepth}{-\maxdimen} % remove section numbering
\usepackage{iftex}
\ifPDFTeX
  \usepackage[T1]{fontenc}
  \usepackage{textcomp} % provide euro and other symbols
\else % if luatex or xetex
  \usepackage{unicode-math} % this also loads fontspec
  \defaultfontfeatures{Scale=MatchLowercase}
  \defaultfontfeatures[\rmfamily]{Ligatures=TeX,Scale=1}
\fi
\usepackage{lmodern}
\ifPDFTeX\else
  % xetex/luatex font selection
\fi
% Use upquote if available, for straight quotes in verbatim environments
\IfFileExists{upquote.sty}{\usepackage{upquote}}{}
\IfFileExists{microtype.sty}{% use microtype if available
  \usepackage[]{microtype}
  \UseMicrotypeSet[protrusion]{basicmath} % disable protrusion for tt fonts
}{}
\makeatletter
\@ifundefined{KOMAClassName}{% if non-KOMA class
  \IfFileExists{parskip.sty}{%
    \usepackage{parskip}
  }{% else
    \setlength{\parindent}{0pt}
    \setlength{\parskip}{6pt plus 2pt minus 1pt}}
}{% if KOMA class
  \KOMAoptions{parskip=half}}
\makeatother
\usepackage{longtable,booktabs,array}
\usepackage{calc} % for calculating minipage widths
% Correct order of tables after \paragraph or \subparagraph
\usepackage{etoolbox}
\makeatletter
\patchcmd\longtable{\par}{\if@noskipsec\mbox{}\fi\par}{}{}
\makeatother
% Allow footnotes in longtable head/foot
\IfFileExists{footnotehyper.sty}{\usepackage{footnotehyper}}{\usepackage{footnote}}
\makesavenoteenv{longtable}
\setlength{\emergencystretch}{3em} % prevent overfull lines
\providecommand{\tightlist}{%
  \setlength{\itemsep}{0pt}\setlength{\parskip}{0pt}}
\usepackage{bookmark}
\IfFileExists{xurl.sty}{\usepackage{xurl}}{} % add URL line breaks if available
\urlstyle{same}
\hypersetup{
  hidelinks,
  pdfcreator={LaTeX via pandoc}}

\author{SCX}
\date{2026}

\begin{document}

\section{Adversarial Verification: Conditional Independence of
Bootstrap-Trained Expert Error
Indicators}\label{adversarial-verification-conditional-independence-of-bootstrap-trained-expert-error-indicators}

\textbf{Claim under review}: If M experts are trained on independent
bootstrap samples D\_1, \ldots, D\_M (drawn i.i.d. with replacement from
the dataset), then their error indicators e\_m(x) = 1\{ℓ(f\_m(x), y)
\textgreater{} τ\} are conditionally independent given the test point x.

\textbf{Reference assumptions}: The paper's A1 requires DISJOINT
independent training sets (D\_m ∩ D\_\{m'\} = ∅), which differs
fundamentally from bootstrap. See below.

\begin{center}\rule{0.5\linewidth}{0.5pt}\end{center}

\subsection{Attack 1: Bootstrap samples are NOT independent --- the
claim's premise is
false}\label{attack-1-bootstrap-samples-are-not-independent-the-claims-premise-is-false}

\textbf{Description}: The claim asserts ``D\_1, \ldots, D\_M are
independent by construction (bootstrap).'' This is mathematically
incorrect.

Bootstrap with replacement from the same finite dataset D\_original =
\{z\_1, \ldots, z\_N\} works as follows:

\begin{itemize}
\tightlist
\item
  D\_m = \{Z\_\{I\_\{m,1\}\}, \ldots, Z\_\{I\_\{m,n\}\}\} where each
  index I\_\{m,j\} is drawn uniformly from \{1, \ldots, N\} with
  replacement.
\item
  D\_1 and D\_2 are functions of (D\_original, ξ\_1) and (D\_original,
  ξ\_2) respectively, where ξ\_1, ξ\_2 are independent index sets.
\end{itemize}

For a \textbf{fixed} D\_original, D\_1 and D\_2 are conditionally
independent given D\_original. But if D\_original is itself random
(which is the standard statistical setting --- the training data is a
random sample from population D), then:

P(D\_1 ∈ A, D\_2 ∈ B) = E{[} P(D\_1 ∈ A \textbar{} D\_original) · P(D\_2
∈ B \textbar{} D\_original) {]}

Since P(D\_1 ∈ A \textbar{} D\_original) and P(D\_2 ∈ B \textbar{}
D\_original) are both functions of the same random D\_original, they are
generally correlated. The expectation of the product does not factor:

E{[}U · V{]} ≠ E{[}U{]} · E{[}V{]} when U, V are correlated.

Therefore, D\_1 and D\_2 are \textbf{not unconditionally independent}.
Step 1 of the proof is wrong.

\textbf{Concrete counterexample}: Let D\_original = \{Z\_1, Z\_2\} with
Z\_1, Z\_2 \textasciitilde{} Bernoulli(p) i.i.d. (so N=2). Draw
bootstrap samples of size n=1 (one element each). Then:

\begin{itemize}
\tightlist
\item
  P(D\_1 = \{1\}) = P(Z\_1=1 or Z\_2=1) = 1 - (1-p)\^{}2 = 2p - p\^{}2
\item
  P(D\_1 = \{1\}, D\_2 = \{1\}) = P(∃ i,j: Z\_i=1 and Z\_j=1) = 1 -
  P(all Z's = 0) - 2·P(exactly one Z=1, other draws that one) \ldots{}
  more directly:
\end{itemize}

P(D\_1=\{1\}, D\_2=\{1\}) = P(at least one Z=1)² ? No.~Let me compute:

P(D\_1=\{1\} \textbar{} D\_original) = 1\{Z\_1=1 or Z\_2=1\}
P(D\_2=\{1\} \textbar{} D\_original) = 1\{Z\_1=1 or Z\_2=1\}

So P(D\_1=\{1\}, D\_2=\{1\}) = E{[}1\{Z\_1=1 or Z\_2=1\}²{]} = P(Z\_1=1
or Z\_2=1) = 2p - p²

And P(D\_1=\{1\})·P(D\_2=\{1\}) = (2p - p²)²

These are equal only when 2p - p² ∈ \{0, 1\}, i.e., p ∈ \{0, 1\}. For
any non-degenerate p ∈ (0,1): 2p - p² ≠ (2p - p²)² because 2p - p² ∈
(0,1).

Thus D\_1 and D\_2 are \textbf{not independent}. They are equal with
positive probability due to both depending on the same D\_original.

\textbf{Verdict: FAIL} --- The very first step of the proof is false.

\begin{center}\rule{0.5\linewidth}{0.5pt}\end{center}

\subsection{Attack 2: The proof sketch ignores the shared label
y}\label{attack-2-the-proof-sketch-ignores-the-shared-label-y}

\textbf{Description}: The proof states ``e\_m(x) depends only on
f\_m(x).'' This is false. The error indicator is:

e\_m(x, y) = 1\{ℓ(f\_m(x), y) \textgreater{} τ\}

It depends on \textbf{both} f\_m(x) and the observed label y. Since y is
a single random variable shared across all M experts, it creates a
common source of randomness. Even if f\_m(x) are independent, the
indicators e\_1, \ldots, e\_M all depend on the same y, so they are not
independent unless y is degenerate.

\textbf{Resolution}: The claim can be salvaged if y is
\textbf{independent} of the vector (f\_1(x), \ldots, f\_M(x)). In the
paper's setting, A4 states that the label noise is independent of all
D\_m and x. Under the paper's A1 (disjoint datasets), combined with a
proper train-test split (x not in any D\_m), we have:

\begin{itemize}
\tightlist
\item
  f\_m(x) depends only on D\_m
\item
  y is independent of \{D\_1, \ldots, D\_M\} (by train-test split + A4)
\item
  Therefore, y is independent of (f\_1(x), \ldots, f\_M(x))
\end{itemize}

In that case, e\_m(x,y) = g(f\_m(x), y) where g(a,b) = 1\{ℓ(a,b)
\textgreater{} τ\}. Since (f\_1(x), \ldots, f\_M(x), y) are independent
random variables, and e\_m depends only on (f\_m(x), y), the indicators
e\_1, \ldots, e\_M are \textbf{not} fully independent of each other
(they share y), but they ARE conditionally independent given y (or given
any σ-algebra that fixes y). However, the claim says ``conditional
independence given x'' not ``given x and y.''

\textbf{More precise analysis}: Define e\_m = 1\{ℓ(f\_m(x), y)
\textgreater{} τ\}. Are e\_1 and e\_2 independent? Let's check:

P(e\_1 = 1, e\_2 = 1 \textbar{} x) = E{[} P(e\_1=1, e\_2=1 \textbar{} x,
f\_1, f\_2, y) \textbar{} x {]} = E{[} 1\{ℓ(f\_1(x), y) \textgreater{}
τ\} · 1\{ℓ(f\_2(x), y) \textgreater{} τ\} \textbar{} x {]}

If y is independent of (f\_1, f\_2) given x, then: = ∫ P(ℓ(f\_1(x), y)
\textgreater{} τ, ℓ(f\_2(x), y) \textgreater{} τ \textbar{} x, f\_1,
f\_2) dP(f\_1, f\_2 \textbar{} x)

The inner probability involves the same y, so it does not factor as
P(ℓ(f\_1(x), y) \textgreater{} τ \textbar{} f\_1) · P(ℓ(f\_2(x), y)
\textgreater{} τ \textbar{} f\_2). The indicators are \textbf{correlated
through y}.

However, if we condition on y as well (i.e., ``given x and y''), then:

P(e\_1=1, e\_2=1 \textbar{} x, y) = P(ℓ(f\_1(x), y) \textgreater{} τ
\textbar{} x, y) · P(ℓ(f\_2(x), y) \textgreater{} τ \textbar{} x, y)

which factors because f\_1(x) and f\_2(x) are independent given x, and y
is now fixed by conditioning.

\textbf{The proof's error}: The proof says ``e\_m(x) depends only on
f\_m(x)'' --- this is simply wrong. It depends on f\_m(x) AND y. The
author appears to have forgotten that e\_m takes two arguments, not one.

\textbf{Verdict: FAIL} --- The proof sketch contains a concrete
mathematical error. The claim IS true if we condition on both x and y,
but the proof as written is incorrect.

\begin{center}\rule{0.5\linewidth}{0.5pt}\end{center}

\subsection{Attack 3: Fatal mismatch between the claim and the paper's
A1}\label{attack-3-fatal-mismatch-between-the-claim-and-the-papers-a1}

\textbf{Description}: The paper's A1 states:

\[D_m \sim \mathcal{D}^{n_m}, \quad D_m \cap D_{m'} = \varnothing, \quad D_m \perp D_{m'} \text{ for } m \neq m'\]

This requires: 1. \textbf{Disjoint} training sets (no overlap) 2. Fresh
i.i.d. draws from the population distribution

The claim substitutes ``bootstrap samples drawn i.i.d. with replacement
from the dataset,'' which is: 1. \textbf{Overlapping} (bootstrap
resamples from the same data) 2. Drawn from a \textbf{single fixed
dataset}, not from the population distribution

These are fundamentally different sampling regimes:

\begin{longtable}[]{@{}lll@{}}
\toprule\noalign{}
Property & Paper (A1) & Claim \\
\midrule\noalign{}
\endhead
\bottomrule\noalign{}
\endlastfoot
Overlap? & Never (disjoint) & Frequent (with replacement) \\
Source & Population D & Fixed dataset \\
Independence & Unconditional & Conditional on dataset only \\
\# of distinct samples & n\_1 + \ldots{} + n\_M & At most N (original
size) \\
\end{longtable}

The paper's Lemma 1, Lemma 2, and main theorem were constructed on the
assumption of disjoint training sets. Swapping in bootstrap invalidates
those proofs. For example, Lemma 1 computes:

E{[}C(x) \textbar{} clean, x{]} = (1/M) Σ\_m P(ℓ(f\_m(x), y)
\textgreater{} τ \textbar{} clean, x)

Under bootstrap, f\_m are \textbf{not independent}, so the average of
e\_m has different concentration properties than what the paper's
Chernoff bounds assume.

\textbf{Verdict: FAIL} --- The claim investigates a different setting
than the paper's. The claim's premise (bootstrap) contradicts the
paper's A1 (disjoint).

\begin{center}\rule{0.5\linewidth}{0.5pt}\end{center}

\subsection{Attack 4: Non-deterministic training and hidden shared
randomness}\label{attack-4-non-deterministic-training-and-hidden-shared-randomness}

\textbf{Description}: The claim assumes f\_m = A(D\_m) where A is
deterministic. In practice:

\begin{itemize}
\tightlist
\item
  Deep learning uses stochastic gradient descent with random minibatch
  order
\item
  Random initialization
\item
  Data augmentation randomness (random crops, flips)
\item
  GPU non-determinism (atomic operation ordering)
\item
  Dropout
\end{itemize}

If A is not deterministic, then f\_m = A(D\_m, ξ\_m) where ξ\_m captures
the training noise. The claim's step 2 (``f\_m depends only on D\_m'')
becomes false.

\textbf{However}: This is \textbf{fixable} if we assume the ξ\_m are
independent across experts and independent of D\_m. Then the argument
extends: f\_m depends on (D\_m, ξ\_m), all of which are independent
across m. The claim would then hold conditional on x and y (assuming the
ξ\_m and D\_m are independent of (x, y)).

\textbf{Edge case}: If all experts use the same random seed, ξ\_m = ξ
for all m, creating shared randomness across experts. In this case,
f\_1, \ldots, f\_M share the same ξ, breaking conditional independence.
This is a real practical concern.

\textbf{Edge case}: GPU non-determinism is not ``randomness'' in the
usual sense --- it's deterministic chaos from parallel floating-point
operations. It cannot be modeled as independent ξ\_m without explicit
intervention (e.g., forcing separate CUDA streams).

\textbf{Verdict: PASS} (with caveats) --- In theory, the randomness can
be made explicit and independent. In practice, shared seeds or GPU
non-determinism can create subtle dependencies. This is not fatal to the
mathematical claim if the independence of ξ\_m is explicitly assumed.

\begin{center}\rule{0.5\linewidth}{0.5pt}\end{center}

\subsection{Attack 5: ``Independent random functions'' does not
automatically imply pointwise independence for random test
points}\label{attack-5-independent-random-functions-does-not-automatically-imply-pointwise-independence-for-random-test-points}

\textbf{Description}: The claim says ``f\_m are independent random
functions, therefore f\_m(x) are independent for fixed x.'' This is
mathematically correct for a \textbf{deterministic} x. However:

\begin{itemize}
\tightlist
\item
  The paper's main results involve expectations over the test
  distribution, not a single fixed x.
\item
  If x is random (X), the statement ``e\_1(X), \ldots, e\_M(X) are
  independent given X'' is a statement about regular conditional
  probabilities.
\item
  For continuous X, the event \{X = x\} has measure zero, requiring
  careful use of regular conditional probabilities.
\end{itemize}

\textbf{This is technically standard} --- the almost-sure definition of
conditional independence works. So this is not a genuine flaw, merely a
technical subtlety.

\textbf{Verdict: PASS} --- Standard measure-theoretic conditioning
handles this.

\begin{center}\rule{0.5\linewidth}{0.5pt}\end{center}

\subsection{Attack 6: In the bootstrap setting, overlapping training
data creates dependence beyond
D\_original}\label{attack-6-in-the-bootstrap-setting-overlapping-training-data-creates-dependence-beyond-d_original}

\textbf{Description}: Even when D\_original is treated as fixed
(conditional inference), bootstrap samples are conditionally independent
given D\_original. However, there is a deeper issue:

Two bootstrap samples D\_1 and D\_2 can share specific training
examples. For instance, if D\_original contains a unique outlier z* (an
image with a rare artifact), and both D\_1 and D\_2 include z\emph{,
both f\_1 and f\_2 will have seen this outlier. Their predictions on
test points similar to z} will be correlated in a way that is NOT
captured by the independence assumption --- they share information about
z*.

Formally: P(f\_1(x) = c \textbar{} D\_original) and P(f\_2(x) = c
\textbar{} D\_original) are both influenced by the same D\_original. The
conditional independence given D\_original is fine, but the
unconditional distribution of f\_1 and f\_2 (integrating over
D\_original) reflects the fact that both experts ``see'' the same data
generating process.

For the paper's concentration bounds, what matters is:

P(\textbar C(x) - E{[}C(x){]}\textbar{} \textgreater{} ε \textbar{} x)

If we use the conditional independence given D\_original, we get:

P(\textbar C(x) - E{[}C(x){]}\textbar{} \textgreater{} ε \textbar{} x,
D\_original) ≤ 2 exp(-2Mε²)

But then taking expectations over D\_original:

P(\textbar C(x) - E{[}C(x){]}\textbar{} \textgreater{} ε \textbar{} x) ≤
E{[}2 exp(-2Mε²) \textbar{} x{]} = 2 exp(-2Mε²)

Wait --- this works because the bound is constant! In this case, the
unconditional bound is the same. But if the bound depends on D\_original
(e.g., if the mean μ varies with D\_original), then:

P(\textbar C(x) - E{[}C(x){]}\textbar{} \textgreater{} ε \textbar{} x) =
E{[}P(\textbar C(x) - E{[}C(x){]}\textbar{} \textgreater{} ε \textbar{}
x, D\_original){]} ≤ E{[}2 exp(-2Mε²){]} = 2 exp(-2Mε²) ← still OK
because bound is constant

Actually, the issue is more subtle for the Chernoff bound used in the
paper. The bound depends on the mean μ\_s, which is an unconditional
expectation. If we only have conditional independence given D\_original,
the Chernoff bound would involve the conditional mean given D\_original,
which is random.

\textbf{This attack is a variant of Attack 1} and is ultimately fatal to
the bootstrap interpretation.

\textbf{Verdict: FAIL} (same root cause as Attack 1)

\begin{center}\rule{0.5\linewidth}{0.5pt}\end{center}

\subsection{Attack 7: The label y might come from the same distribution
as the training data, creating subtle
dependence}\label{attack-7-the-label-y-might-come-from-the-same-distribution-as-the-training-data-creating-subtle-dependence}

\textbf{Description}: In the paper's label noise model (A4), the
observed label y is:

y = y* (true label) with probability 1-η y = Uniform(Y ~\{y*\}) with
probability η

The label noise is independent of D\_m. But the true label y* =
f\emph{(x) depends on x, which is the test point. If x is from the test
set (properly separated), then y} is independent of D\_m, and the claim
holds.

\textbf{But what if the claim's ``bootstrap'' setting includes the
possibility that x comes from the training data?} In bootstrap
validation (like out-of-bag estimation), some test points may have been
seen in some bootstrap samples. If x ∈ D\_1 (overlap between train and
``test''), then:

\begin{itemize}
\tightlist
\item
  f\_1(x) is trained on x itself, creating a data leakage
\item
  y is not independent of D\_1 (since the pair (x, y) might be in D\_1)
\item
  e\_1(x) and e\_2(x) now have a complex dependence structure through
  the shared (x, y) in D\_1
\end{itemize}

This breaks the conditional independence argument entirely. For a proper
test point (not in any training set), this doesn't apply. But the claim
doesn't specify train-test separation.

\textbf{Verdict: FAIL} --- The claim uses ``test point x'' language but
doesn't ensure x is not in any bootstrap sample. For proper evaluation
(new unseen test data), this can be assumed away, but the bootstrap
framing makes it ambiguous.

\begin{center}\rule{0.5\linewidth}{0.5pt}\end{center}

\subsection{Final Verdict}\label{final-verdict}

\textbf{The claim IS NOT rigorous.} There are 4 significant issues, 3 of
which are fatal:

\begin{longtable}[]{@{}
  >{\raggedright\arraybackslash}p{(\linewidth - 6\tabcolsep) * \real{0.1000}}
  >{\raggedright\arraybackslash}p{(\linewidth - 6\tabcolsep) * \real{0.2667}}
  >{\raggedright\arraybackslash}p{(\linewidth - 6\tabcolsep) * \real{0.3333}}
  >{\raggedright\arraybackslash}p{(\linewidth - 6\tabcolsep) * \real{0.3000}}@{}}
\toprule\noalign{}
\begin{minipage}[b]{\linewidth}\raggedright
\#
\end{minipage} & \begin{minipage}[b]{\linewidth}\raggedright
Attack
\end{minipage} & \begin{minipage}[b]{\linewidth}\raggedright
Severity
\end{minipage} & \begin{minipage}[b]{\linewidth}\raggedright
Verdict
\end{minipage} \\
\midrule\noalign{}
\endhead
\bottomrule\noalign{}
\endlastfoot
1 & Bootstrap samples are not independent & \textbf{Fatal} --- step 1 of
proof is false & FAIL \\
2 & Proof ignores the shared label y & \textbf{Fatal} --- proof contains
mathematical error & FAIL \\
3 & Mismatch with paper's A1 (disjoint vs bootstrap) & \textbf{Fatal}
--- different sampling regime & FAIL \\
6 & Bootstrap overlap creates dependence through D\_original &
\textbf{Fatal} --- variant of Attack 1 & FAIL \\
7 & Training-test leakage in bootstrap & \textbf{Fatal} --- not
specified & FAIL \\
4 & Non-deterministic training & Caveat, fixable & PASS \\
5 & Formal conditioning issues & Standard technicality & PASS \\
\end{longtable}

\textbf{Bottom line}: The claim suffers from a fatal contradiction in
its premises (``bootstrap'' and ``independent'' are incompatible), a
concrete mathematical error in its proof (ignoring the shared label y),
and a mismatch with the paper's A1 assumption. Under the paper's actual
assumption (disjoint datasets), and with a corrected proof that properly
accounts for y, the conditional independence claim can be made rigorous
--- but the claim as stated and its proof sketch are not correct.

\end{document}
